
\section{azapy package}\label{AZAPY}

\azapy\ is an open-source Python package for portfolio optimizations.
It is released under GNU General Public License v3.0 (GPL-3).
Besides the set of algorithms presented in this paper, \azapy\ implements several other
portfolio strategies. It is a work in progress free for everyone to use and explore.

The package technical documentation is at  \azapydoc\ and its source code at
\azapycode.
It can be installed in a Python environment by calling
\begin{lstlisting}[language=Python]
    pip install azapy
\end{lstlisting}
The \azapy\ development has two major goals:
\begin{itemize}
  \item provide a technical environment to study and compare various portfolio optimizations (in-sample and out-of-sample),
  \item provide an analytic library to support the maintenance of a portfolio rebalancing schedule for a retail type of investor.
\end{itemize}

The library has its own module to collect and maintain historical market data. Several market data providers
are supported. Some of them are free while other require subscriptions.
The list of current supported providers is present in the package documentation at \azapydoc.

The optimization algorithms are using open-source packages for LP, QP and SCOP:
\ecos~\cite{ecos}, \highs~\cite{highs-ds} through \linprog, \cvxopt, and \glpk.

We would like to express our thanks to all developers who were generous to give us unselfish access to their work and imagination.

All the optimization algorithms presented in this paper are accompanied by a simple illustrative Python script.
A set of more complex Python and Jupyter Notebook scripts are present at \azapycode\ under directories {\tt scripts} and {\tt jpy-scripts}, respectively.
The idea is to use them as examples for calling various \azapy\ functionalities as well as sources of inspiration in your work.

If you like to comment or contribute to this project, feel free to contact us (details are at PyPi page \azapy).

\BackToContents 